% Preámbulo --------------------------------------------------------------------
\documentclass[a4paper, titlepage]{article}

\usepackage[T1]{fontenc}
\usepackage[sfdefault=true]{inter}
\usepackage{amsmath}
\usepackage{graphicx}
\usepackage[export]{adjustbox}
\usepackage{array}
\usepackage{booktabs}
\usepackage{multirow}
\usepackage{framed}
\usepackage[spanish]{babel}
\usepackage[style=english]{csquotes}
\usepackage[style=apa]{biblatex}
\addbibresource{referencias.bib}
\usepackage[dvipsnames]{xcolor}
\usepackage{hyperref}
\hypersetup{colorlinks=true, linkcolor=RoyalBlue, urlcolor=RoyalBlue, citecolor=RoyalBlue}

\title{Trabajo Final Integrador \\[1em]%
  \large Carrera de Especialización en Estadística \\%
  Facultad de Ciencias Exactas y Naturales \\%
  Universidad de Buenos Aires}
\author{Fernando Gallegos Piderit}
\date{\today}

\begin{document} % ------------------------------------------------------------

\maketitle

\tableofcontents

\section{Introducción}

% Objetivo del TP, obtener dos modelos de clasificación que permitan identificar
% tumores malignos. Breve descripción de los datos con los que se cuenta.
% Descripción de la estructura del informe. ¿Mención al software, paquetes de R
% etc., utilizado?

% Sección: Introducción

El objetivo del trabajo práctico fue construir dos clasificadores que, a partir
de cierta información acerca de las características de las células que componen
un tumor, permitan identificar si este es maligno o benigno con un grado
aceptable de error bajo algún criterio establecido. Para ello se contó con una
muestra de observaciones de tumores de los que se registraron tres medidas
resumen de diez características de una cantidad desconocida de sus células
mediante imágenes microscópicas \parencite{williamwolberg_1993_breast}. Los
clasificadores obtenidos se seleccionaron de un total de seis que fueron
entrenados y evaluados sobre los datos disponibles según su desempeño en
ciertas métricas que reflejan los criterios antes mencionados. Como es
habitual, el entrenamiento y la evaluación de los modelos se hicieron sobre un
conjunto de entrenamiento y otro de \textit{testeo} respectivamente. Los
modelos evaluados fueron los siguientes: regresión logística, regresión
logísitca con penalización LASSO, bosques aleatorios, análisis discriminante
lineal y $k$ vecinos más cercanos. De estos, los dos que mostraron el mejor
desempeño fueron los bosques aleatorios y la regresión logística con
penalización LASSO con las variables originales normalizadas en ambos casos.

El informe se estructura de la siguiente manera. En la Sección \ref{sec:aed} se
da una descripción más detallada de los datos con los que se entrenaron los
modelos de clasificación y se resaltan algunos aspectos identificados que
tuvieron impacto sobre la estrategia de modelado. En la Sección
\ref{sec:estmod} se describe esta estrategia, en particular, algunas decisiones
generales de preprocesamiento de los datos, el método de validación, la
elección de las métricas de evaluación y su justificación y los modelos
específicos a implementar. En la Sección \ref{sec:res} se presentan los
resultados de la evaluación de todos los modelos implementados sobre los datos
de testeo y se detallan más específicamente los dos que mejor desempeño
tuvieron junto con el procedimiento completo de modelado. Finalmente, en la
Sección \ref{sec:con} se discuten brevemente los resultados y se resaltan los
puntos que resultaron más problemáticos en el desarrollo del trabajo práctico.


\section{Análisis exploratorio}\label{sec:aed}

% Descripción más detallada de las variables del dataset, tipos, casos perdidos,
% etc. Medidas resumen de la variable de respuesta y de las demás. Histogramas y
% diagramas de dispersión de algunas ¿cuántas y cuáles? variables. ¿Quizás las
% más "representativas"? Mencionar la alta correlación entre las tres medidas
% resumen, mean, sd y worst, eso fundamenta el PCA.

% Sección: Análisis exploratorio de datos

Los datos completos cuentan con 569 observaciones y 31 variables, de las cuales
una, \verb|target|, es la variable de respuesta que nos indica si el tumor al
que corresponde cada observación es maligno (\verb|M|) o benigno (\verb|B|), y
las demás son variables numéricas predictoras que resumen la información
obtenida de las células de dichos tumores. Como se mencionó en la introducción,
estas variables corresponden a tres medidas resumen (la media, el desvío estándar
y el peor caso) de diez características celulares listadas a continuación:

\begin{itemize}
  \item \verb|radius|: El radio de las células
  \item \verb|texture|: Variación en la intensidad del núcleo
  \item \verb|perimeter|: Perímetro del núcleo
  \item \verb|area|: Área del núcleo
  \item \verb|smoothness|: Variación local del radio de las células
  \item \verb|compactness|: Qué tan compacta o extendida es la célula
  \item \verb|concavity|: Profundidad de las partes cóncavas del contorno de las células
  \item \verb|concave|: Número de secciones cóncavas del contorno de las células
  \item \verb|symmetry|: Qué tan simétrica es la forma de las células
  \item \verb|fractal_dimension|: Complejidad del borde de las células
\end{itemize}

Antes de cualquier partición de los datos, observamos que la proporción de
observaciones de tumores malignos es 0.3726 aproximadamente (212
observaciones); es decir, los datos están ligeramente desbalanceados. Sin
embargo, dado que el desbalance no es tan marcado, se optó por no usar estratos
en la partición en datos de entrenamiento y testeo para minimizar las
decisiones tomadas sobre la totalidad de los datos. Se retuvo aproximadamente
el 30\% de las observaciones (171) para la evaluación de los modelos y el resto
se utilizó para el análisis exploratorio, el entrenamiento y la validación.

En los datos de entrenamiento la proporción de tumores malignos es 0.3643
aproximadamente. Afortunadamente no hay datos perdidos en ninguna de las
variables. En el Cuadro \ref{tab:resumen} se muestran algunas medidas resumen
de las variables predictoras, agrupadas según la medida resumen (valga la
redundancia) a la que correspondan. Las variables predictoras tienen unidades sumamente
dispares; por ejemplo, \verb|concave_mean| tiene valores entre 0 y 0.19
secciones cóncavas mientras que \verb|area_mean| los tiene entre 143 y 2501
unidades de área. Esta información es relevante ya que querríamos evitar que
algunas características tengan mayor impacto en el entrenamiento de los modelos
que otras. Todas las variables tienen rangos positivos y las posiciones de sus
medias y medianas sugieren que la mayoría son asimétricas con un sesgo hacia
los valores más altos.

Las Figuras \ref{fig:area}-\ref{fig:smooth} y
\ref{fig:radius}-\ref{fig:fractal}\footnote{La mayoría de los gráficos
exploratorios se presentan en el Anexo \ref{sec:anexo2} ya que son numerosos y
ocupan bastante espacio} ofrecen visualizaciones de las distribuciones de las
variables del conjunto de datos. Cada figura corresponde a una característica
de las células; la fila superior muestra histogramas de densidad según el valor
de \verb|target| para cada medida resumen y la fila inferior muestra un
diagrama de dispersión por cada combinación de dos de esas medidas con los
puntos coloreados según los valores de \verb|target|.

% Gráficos area
\begin{figure}[h]
\includegraphics[width=\textwidth, center]{figuras/distribuciones\_area.pdf}
\caption{Histogramas de densidad de las variables asociadas a \textit{area} según \textit{target} (fila superior) y diagramas de dispersión de las mismas variables de a pares (fila inferior)}
\label{fig:area}
\end{figure}

Si nos enfocamos en los histogramas, se puede observar que, efectivamente,
algunas variables tienen marcados sesgos hacia los valores más altos, en
particular las que corresponden a los desvíos estándar de las características
de los tumores. Estas variables también tienden a presentar valores extremos;
esto se puede apreciar en la Figura \ref{fig:area}, por ejemplo, en el desvío
estándar del área de las células. También se observa que en casi todas las
variables los tumores malignos tienden a tener valores más altos. Las
excepciones son las variables relacionadas con la dimensión fractal del
contorno de las células y el desvío estándar de la simetría, la suavidad y la
textura. En algunos casos, como el área promedio de las células, las
distribuciones condicionales a la calidad de benigno o maligno del tumor son
más simétricas que la distribución general de la variable.

Un aspecto relevante de las variables es su dependencia; es de esperar que las
que corresponden a las mismas características de las células estén relacionadas
entre sí. Si examinamos los diagramas de dispersión de las Figuras
\ref{fig:area}-\ref{fig:smooth} y \ref{fig:radius}-\ref{fig:fractal} vemos que,
efectivamente, tienden a estar fuertemente asociadas con un patrón bastante
lineal. La correlación es más fuerte entre la media y el peor caso, con el
coeficiente $\rho$ de pearson yendo de 0.688 para la simetría hasta 0.9718 para
el perimetro. Es más débil entre la media y el desvío, yendo de 0.3505 para la
suavidad hasta 0.8077 para el área, y entre el desvío y el peor caso, de 0.4035
para la suavidad hasta 0.86 para el área. 

\begin{figure}[h]
\includegraphics[width=\textwidth, center]{figuras/distribuciones\_smoothness.pdf}
\caption{Histogramas de densidad de las variables asociadas a \textit{smoothness} según \textit{target} (fila superior) y diagramas de dispersión de las mismas variables de a pares (fila inferior)}
\label{fig:smooth}
\end{figure}

Si nos enfocamos en la condición de maligno o benigno de los tumores en los
diagramas de dispersión, vemos que, a diferencia de lo que se ve en los
histogramas, hay fronteras bastante nítidas entre ambos tipos en algunas de las
características. Esto es muy claro en las variables asociadas al radio (Figura
\ref{fig:radius}), el perímetro (Figura \ref{fig:perimeter}) y el área (Figura
\ref{fig:area}). En otros casos, como en las variables de suavidad (Figura
\ref{fig:smooth}), no parece haber una frontera muy clara. Nada de esto
descarta que sí exista una frontera nítida en dimensiones más altas; las
visualizaciones que se hicieron no van más allá de las dos dimensiones y,
teniendo 30 variables predictoras, es probable que no estemos notando patrones
importantes.


\section{Estrategia de modelado}\label{sec:estmod}

% Sección: Estrategia de modelado
%
% Mencionar, no sé en qué orden: partición en train y test; preprocesamiento de
% los datos, normalización, pca, etc.; existencia de hiperparámetros y el método
% utilizado para seleccionarlos, CV con 5 folds; ¿capaz aquí sí el software,
% paquetes de R, etc.?; procedimiento general con cada modelo.

Como se mencionó en la sección anterior, los datos originales se dividieron en
un conjunto de entrenamiento, con aproximadamente un 70\% de las observaciones,
y uno de testeo o evaluación con el resto. El primero se utilizó para entrenar
y validar los modelos, y el segundo para calcular las métricas en base a las
que se eligieron los dos mejores.

En términos generales, la estrategia de modelado debió tomar en cuenta los
siguientes problemas. En primer lugar, el análisis exploratorio mostró que los
variables predictoras tienen unidades dispares y pueden llegar a presentar
correlaciones lineales muy altas. Esto llevó a la decisión de normalizar las
variables y de probar estrategias de reducción de la dimensión o regularización
en algunos casos. En segundo lugar, varios de los modelos propuestos
(enumerados en la Subsección \ref{sec:modelos}) cuentan con uno o más
hiperparámetros cuyos valores deben ser establecidos antes de la evaluación en
los datos de testeo. En algunos casos se eligieron determinados valores en
función de lo observado en el proceso de modelado y del costo computacional. En
los casos en los que no, dichos valores se eligieron por validación cruzada de
cinco folds buscando sobre grillas de valores definidas de antemano. Esto
introdujo el problema de seleccionar una métrica en base a la cual elegir los
valores más adecuados de los hiperparámetros. Por motivos que se detallan en la
Subsección \ref{sec:metricas} se eligió una métrica distinta a la de evaluación
de los modelos. Por último, se debió establecer un umbral de predicción que
determine qué casos se clasifican como malignos y cuáles como benignos.

El trabajo se realizó íntegramente en R 4.5.2 y RStudio 2025.09.2
\parencite{rcoreteam_2025_language,positteam_2025_rstudio}, fundamentalmente
con los conjuntos de paquetes de \verb|{tidyverse}| y \verb|{tidymodels}|
\parencite{wickham_etal_2019_welcome,kuhn_wickham_2020_tidymodels}. Estos
ofrecen una interfaz sumamente consistente para gran parte del proceso de
modelado, lo cual simplificó mucho el lidiar con seis modelos distintos y todos
los pasos que implica su construcción. En cuanto a los modelos en sí,
\verb|{tidymodels}| hace uso de otros paquetes para implementarlos. En este
caso se usaron los paquetes \verb|{ranger}|, \verb|{glmnet}|, \verb|{kknn}| y
\verb|{discrim}|
\parencite{wright_ziegler_2017_ranger,friedman_etal_2010_regularization,tay_etal_2023_elastic,schliep_hechenbichler_2025_kknn}.

\subsection{Métricas de evaluación y validación}\label{sec:metricas}

% Aclarar que hay dos métricas usadas, área bajo la curva Precision-Recall para
% elegir hiperparámetros, y puntaje $F_\beta$ para evaluar el desempeño en los
% datos de testeo. ¿Otras métricas como precision y recall? Justificar $F_\beta$
% con riesgo de no detectar cáncer vs. riesgo de diagnosticar erróneamente a
% alguien. Mencionar problema con parámetros $\beta$ y el umbral de decisión.
% Explicar cómo se elige el umbral en los modelos. ¿$\beta = 1$?¿Cómo justificar
% otra ponderación?

La elección de los dos mejores modelos se hizo en base al puntaje $F_\beta$
(definido en los párrafos siguientes) con $\beta = 2$ y calculado sobre el
conjunto de testeo. Sin embargo, también se calculó la \textit{accuracy}, la
precisión, el \textit{recall} y el área bajo la curva \textit{precision-recall}
para poder interpretar mejor los resultados. La selección del puntaje $F_2$
como métrica principal se basó en la consideración de las consecuencias de los
falsos positivos y los falsos negativos en la identificación de tumores
malignos. Durante todo el ejercicio se consideró como clase positiva a los
tumores malignos.

\begin{table}[h]
\begin{tabular*}{\textwidth}[t]{@{\extracolsep{\fill}}lcc}
\toprule
\multirow{2}{3em}{Realidad}  & \multicolumn{2}{c}{Predicción} \\
            \cmidrule{2-3}
          & Maligno                 & Benigno \\
\midrule
Maligno   & Verdadero Positivo (VP) & Falso Negativo (FN) \\
Benigno   & Falso Positivo (FP)     & Verdadero Negativo (VN) \\
\bottomrule
\end{tabular*}
\caption{Matriz de confusión.}
\label{tab:confusion}
\end{table}

El Cuadro \ref{tab:confusion} muestra la matriz de confusión que corresponde al
problema en cuestión. Un falso negativo, en este escenario, es el peor de los
resultados ya que implica que un tumor maligno no fue detectado. Si bien el
diagnóstico definitivo del cáncer no se obtiene mediante un único método, su
detección temprana es fundamental para aumentar las chances de supervivencia de
las pacientes y cualquier retraso pone en riesgo sus vidas
\parencite{saunders_jassal_2009_breast}. Un falso positivo, si bien no alcanza
el mismo nivel de gravedad, sigue sin ser un buen resultado; significa que se
identificó como maligno un tumor benigno, lo cual puede llevar a que la persona
deba someterse a estudios invasivos como biopsias o a otros tratamientos. Como
estamos lidiando con tumores mamarios, esto puede significar cirugías,
radioterapia o quimioterapia \parencite{saunders_jassal_2009_breast},
dependiendo de la gravedad del error. Como es razonable en cualquier problema
de clasificación, el objetivo fue minimizar los casos mal clasificados, salvo
que, en este caso, hay un tipo de error que es prioritario evitar por sobre el
otro.

Dos métricas habituales que sirven en este caso son la precisión y el
\textit{recall}. El \textit{recall} o \textit{true positve rate} mide la
proporción de los casos verdaderamente positivos que fueron predichos como
positivos por el modelo:

\begin{equation}
\text{Recall} = \frac{\text{VP}}{\text{VP} + \text{FN}}
\end{equation}

Por su parte, la precisión o \textit{positive predictive value} mide la
proporción de los predichos como positivos por el modelo que realmente son
positivos:

\begin{equation}
\text{Precision} = \frac{\text{VP}}{\text{VP} + \text{FP}}
\end{equation}

Ambas métricas se aproximan a 0 a medida que aumentan los casos mal
clasificados, y a 1 a medida que disminuyen. Sin embargo, existe cierto
\textit{trade-off} entre ambas; si predecimos todos los casos como positivos,
entonces no habrá falsos negativos y el \textit{recall} será 1, pero
inevitablemente habrá falsos positivos, por lo cual la precisión se acercará a
0. Además, puede suceder que la proporción de casos realmente positivos sobre
el total sea muy alto, y que, por lo tanto, un mal modelo logre una alta
precisión simplemente porque hay pocos casos realmente negativos para ser
clasificados como positivos. En ese caso seguramente haya más falsos negativos
y el \textit{recall} será alto.

El puntaje $F_\beta$ ofrece un balance entre ambas métricas. Es su media
armónica ponderada de forma tal que el \textit{recall} pesa $\beta$ veces más
que la precisión:

\begin{equation}
\begin{aligned}
  F_\beta &= \frac{(1 + \beta^2) \text{Precision} \times  \text{Recall}}{\beta^2 \text{Precision} + \text{Recall}} \\
  &= \frac{(1 + \beta^2) \text{VP}}{(1 + \beta^2) \text{VP} + \beta^2 \text{FN} + \text{FP}}
\end{aligned}
\end{equation}

La elección del valor de $\beta$ representó un problema difícil de resolver en
el trabajo. En principio parece evidente que un valor de $\beta = 1$ es
inadecuado ya que las consecuencias de los falsos negativos son más graves que
las de los falsos positivos. Sin embargo, no es claro cómo ir más allá de un
orden en la gravedad de los errores. Un antecedente del campo de la
bioinformática que utiliza un valor de $\beta = 2$ es
\textcite{sanchez-herrero_etal_2024_forecasting}. Sin embargo, otros antecedentes
indican que lo más común es usar $\beta = 1$, y que, de todas formas, el
puntaje $F_\beta$ es inadecuado en casos clínicos
\parencite{assel_vickers_2025_score}. Finalmente se optó por usar un valor de
$\beta = 2$, que es relativamente arbitrario pero mayor a 1.

Las tres medidas presentadas requieren que se establezca un umbral de
predicción que permita convertir las probabilidades de pertenecia a la clase
maligna que estiman los modelos en una clasificación de las observaciones en
dos clases. Este umbral se seleccionó durante la etapa de validación de los
modelos en los datos de entrenamiento (esto se desarrollará en más detalle en
la Sección \ref{sec:detalle}).

Como se mencionó anteriormente, para elegir los hiperparámetros de algunos
modelos se realizó una validación cruzada sobre cinco folds. La métrica que se
buscó maximizar en este caso fue el área bajo la curva
\textit{precision-recall}, no el puntaje $F_2$, ya que la primera no depende
del umbral de decisión y también refleja un balance entre ambas medidas. De lo
contrario, se debería haber incluido el umbral como hiperparámetro a optimizar
durante la validación cruzada, lo cual hubiese complejizado la búsqueda.

\subsection{Modelos a implementar}\label{sec:modelos}

% Enumerar y describir brevemente los modelos a implementar. Mencionar los
% hiperparámetros que tiene cada uno, cuáles van a ser optimizados y cuáles no.
% Justificar los valores utilizados para los que no se optimizan.

Los modelos que se implementaron fueron los siguientes:

\begin{enumerate}

  \item Regresión logística sobre las componentes principales de las variables
    predictoras normalizadas. Se optimizó el número de componentes a utilizar
    mediante validación cruzada. Se decidió hacer una reducción de la dimensión
    dada la alta correlación que se encontró entre las variables predictoras
    durante el análisis exploratorio.
  \item Regresión logística con penalización LASSO sobre las variables
    predictoras originales normalizadas. Se optimizó el hiperparámetro de
    penalización $\lambda$ mediante validación cruzada. Se incluyó este modelo
    ya que la penalización LASSO no sólo permite seleccionar variables
    relevantes, sino que también enfrenta el problema de la correlación entre
    las predictoras limitando el tamaño de los coeficientes.
  \item Regresión logística con penalización LASSO sobre las componentes
    principales de las variables predictoras normalizadas. Al igual que en el
    modelo anterior, se optimizó el hiperparámetro $\lambda$. Se incluyó este
    modelo como método alternativo de selección de componentes principales; el
    primer modelo selecciona las primeras $n$ componentes, mientras que este
    selecciona $m$ componentes que no necesariamente son las primeras.
  \item Bosques aleatorios sobre las variables predictoras originales
    normalizadas. Este modelo cuenta con tres hiperparámetros (al menos en la
    implementación usada): la cantidad de árboles por bosque, la cantidad
    mínima de observaciones por hoja, y la cantidad de variables a utilizar por
    \textit{split} en cada árbol. Durante la validación se notó que la cantidad
    de árboles no afectaba notoriamente el desempeño del modelo, por lo cual se
    eligió usar 100 árboles y se optimizaron los otros dos parámetros (sobre
    grillas reducidas para no aumentar demasiado el costo computacional).
  \item Análisis discriminante lineal sobre las componentes principales de las
    variables predictoras normalizadas. Se optimizó el número de componentes a
    incluir en el modelo. Se utilizaron las componentes principales ya que
    durante el modelado se observó que usando sólo las variables normalizadas
    el modelo tenía un desempeño mucho peor que el del resto.
  \item $K$ vecinos más cercanos sobre las variables predictoras originales
    normalizadas. Se optimizó la cantidad de vecinos.

\end{enumerate}


\section{Resultados}\label{sec:res}

% Sección: Resultados
%
% ¿Valdrá la pena mostrar algunos resultados intermedios, anteriores a la
% evaluación en los datos de testeo? Umbrales de decisión, métricas vs.
% hiperparámetros, etc.

\subsection{Desempeño de los modelos implementados}

% Describir verbalmente el desempeño y mostrar métricas calculadas sobre los
% datos de testeo. ¿Mostrar gráficos, modelos y métricas por modelo? ¿Mostrar
% cuadros de métricas por modelo? Hacer una breve interpretación de los
% resultados en relación con la temática.

% Métricas de los modelos calculadas en test
\begin{table}[h]
\begin{tabular*}{\textwidth}{@{\extracolsep{\fill}}lrrrrr}
\toprule
  Modelo & $F_2$ & \textit{Accuracy} & \textit{Precision} & \textit{Recall} & \textit{AUC (PR)}\\
\midrule
Random Forest & 0.9941 & 0.9883 & 0.9710 & 1.0000 & 0.9993\\
Logit LASSO & 0.9911 & 0.9825 & 0.9571 & 1.0000 & 0.9993\\
KNN & 0.9853 & 0.9708 & 0.9306 & 1.0000 & 1.0000\\
Logit LASSO (PCA) & 0.9851 & 0.9883 & 0.9851 & 0.9851 & 0.9973\\
LDA & 0.9735 & 0.9649 & 0.9296 & 0.9851 & 0.9961\\
Logit & 0.9552 & 0.9649 & 0.9552 & 0.9552 & 0.9891\\
\bottomrule
\end{tabular*}

\caption{Métricas de los modelos en los datos de testeo. Se presentan en orden descendiente según el puntaje $F_2$ obtenido en el conjunto de testeo.}
\label{tab:metricas}
\end{table}

En el Cuadro \ref{tab:metricas} se presentan las métricas obtenidas por cada
modelo sobre los datos de testeo. Los dos modelos que alcanzaron los mayores
valores en $F_2$ fueron el de bosques aleatorios con 0.9941, y la regresión
logística con penalización LASSO sobre las variables originales con 0.9911. En
ambos casos la clasificación fue casi perfecta: los bosques aleatorios lograron
una precisión de 0.971 con sólo dos falsos positivos y un \textit{recall}
perfecto; la regresión logró una precisión de 0.957 con tres falsos positivos
y, también, \textit{recall} perfecto. El resto de los modelos no se quedó
demasiado atrás: el que tuvo el peor desempeño fue la regresión logística sobre
las componentes principales de las variables originales con un puntaje
$F_2$ de 0.955 y seis tumores mal clasificados, tres benignos y tres malignos.

% Matrices de confusión de los modelos elegidos
\begin{table}[h]
\begin{tabular*}{\textwidth}{@{\extracolsep{\fill}}lcccc}
\toprule
\multirow{3}{3em}{Realidad}  & \multicolumn{4}{c}{Predicción} \\
          & \multicolumn{2}{c}{Bosques aleatorios} & \multicolumn{2}{c}{Logística LASSO} \\
            \cmidrule{2-5}
          & Maligno & Benigno  & Maligno & Benigno \\
\midrule
Maligno   & 67      & 0        & 67      & 0 \\
Benigno   & 2       & 102      & 3       & 101 \\
\bottomrule
\end{tabular*}

\caption{Matrices de confusión de los dos mejores modelos.}
\label{tab:matrices}
\end{table}

Es interesante notar el efecto de la ponderación en el puntaje $F_\beta$ sobre
el resultado. La regresión logística con penalización LASSO sobre las
componentes principales de las predictoras tuvo la misma \textit{accuracy} que
el modelo de bosques aleatorios (dos tumores mal clasificados). Sin embargo,
como se puede ver en el Cuadro \ref{tab:matrices}, una de sus predicciones fue
un falso positivo, por lo cual terminó teniendo un puntaje $F_2$ menor al de
los bosques aleatorios.

\subsection{Descripción de los modelos seleccionados}\label{sec:detalle}

% Describir más detalladamente el procedimiento de cada modelo. ¿Quizás aquí van
% los gráficos de resultados intermedios? Mostrar los parámetros, los del pca y
% los de las regresiones si es que fueran seleccionadas. ¿Mencionar lo de
% tidymodels acá?

\subsubsection{Bosques aleatorios}

Para entrenar este modelo se utilizaron las variables originales normalizadas
salvo las asociadas a la dimensión fractal (\verb|fractal_dimension\_mean|,
\verb|..._se| y \verb|..._worst|) y el desvío estándar de la textura
(\verb|texture_se|), la suavidad (\verb|smoothness_se|) y la simetría
(\verb|symmetry_se|), que se excluyeron dado que en los histogramas de densidad
no se vió que ofrecieran capacidad de separar las dos clases de tumores. Como
los árboles aleatorios utilizan una única variable por \textit{split}, no tiene
ventaja incluir aquellas cuya distribución es similar en ambas clases. Por un
motivo similar se decidió no buscar las componentes principales de las
variables; la alta correlación entre las predictoras no debería entrar en juego
si en cada \textit{split} sólo se usa una de ellas.

Para el modelo en sí se utilizó la implementación del paquete \verb|{ranger}| a
través de \verb|{tidymodels}|. Los hiperparámetros de este modelo son los
siguientes: el número de árboles a entrenar, que, tras notar que no producía
cambios notorios en la métrica de validación, se estableció en 100, un número
relativamente bajo para no aumentar innecesariamente el costo computacional; el
número de observaciones mínimo por hoja, que se buscó optimizar; y el número de
variables a considerar en cada split de cada árbol, que también se optimizó. En
este proceso se consideraron los valores de 10, 20 y 30 observaciones mínimas
por hoja, y entre 1 y 10 variables por split, y se buscó sobre una grilla con
todas las combinaciones de estos dos grupos de valores. Como se mencionó antes,
la optimización se hizo mediante validación cruzada de cinco \textit{folds}
maximizando el área bajo la curva \textit{precision-recall}. Los \textit{folds}
se crearon estratificando por \verb|target| para mantener la proporción de
casos malignos dentro de cada uno.

% Métricas de validación de random forest
\begin{figure}[h]
\includegraphics[width=\textwidth]{figuras/rf\_metricas.pdf}
\caption{Métricas de desempeño del modelo de bosques aleatorios según hiperparámetros optimizados. Los gráficos están segmentados según el número mínimo de observaciones por hoja.}
\label{fig:rfmetricas}
\end{figure}

En la Figura \ref{fig:rfmetricas} se puede ver la métrica de validación
calculada según los hiperparámetros junto con la precisión y el \textit{recall}
calculados con un umbral de 0.5. La combinación que rindió la mayor área bajo
la curva fue una variable por split y diez observaciones mínimo por hoja de
cada árbol con un área bajo la curva promedio de 0.985 y error estándar de
0.004 aproximadamente. Algo que se observa es que el área no varía demasiado
según el hiperparámetro, pero sí lo hacen la precisión y el \textit{recall}.

% Acumulada de las probabilidades predichas de random forest
\begin{figure}[h]
  \includegraphics[width=\textwidth]{figuras/rf\_pred\_umbral.pdf}
  \caption{(A) Distribución acumulada de las probabilidades predichas y (B) puntaje $F_2$ según umbral de decisión correspondientes al modelo de bosques aleatorios. La línea punteada vertical indica el umbral de decisión elegido.}
\label{fig:rfpredumb}
\end{figure}

Una vez obtenidos los hiperparámetros optimizados se utilizaron las
predicciones hechas sobre los cinco \textit{folds} para buscar el umbral de
decisión óptimo que maximizase el puntaje $F_2$. Esto se hizo sobre el total de
las predicciones sin distinguir el \textit{fold} al que perteneciesen. Se buscó
sobre una grilla de 99 valores equiespaciados entre 0.01 y 0.99, y se tomó como
óptimo el promedio de los umbrales que alcanzasen el máximo $F_2$ (ya que podía
haber más de un valor que tocase el máximo). En la Figura \ref{fig:rfpredumb}(A) se
muestran las distribuciones acumuladas de las probabilidades predichas de ser
un tumor maligno según la condición real del tumor y el \textit{fold} de
validación. Se ve que el modelo logra una buena separación de los casos, pero
pareciera que en una buena proporción de los casos malignos las probabilidades
predichas son demasiado bajas; en los cinco \textit{folds} aproximadamente el
12.5\% tiene una probabilidad predicha menor a 0.5. La figura
\ref{fig:rfpredumb}(B) muestra el puntaje $F_2$ alcanzado según el umbral de
decisión y \textit{fold}. Se puede ver que el máximo suele alcanzarse en
valores más bajos que 0.5. La línea negra punteada se ubica en 0.29 en el eje
$x$ e indica el mejor umbral encontrado, donde se alcanza un $F_2$ de 0.9435.

En el Cuadro \ref{fig:rfespec} del Anexo \ref{sec:anexo1} se resumen las características del modelo final de forma más sintética.


\subsubsection{Regresión logística con penalización LASSO}

Para este modelo no se realizó una selección de variables ya que la
penalización LASSO tiene como consecuencia que algunos de los coeficientes
estimados de la regresión pueden ser exactamente 0. Se utilizaron todas las
variables predictoras disponibles en el conjunto de datos y se las normalizó
para poder implementar el modelo. El término de penalización LASSO es $\lambda
\sum_{j=1}^p | \beta_j |$, donde $p$ es la cantidad de variables predictoras,
en este caso 30, y $\beta_j$ es el coeficiente correspondiente a la predictora
$j$. Como los coeficientes $\beta_j$ están en la unidad de sus respectivas
variables, si estas no se normalizaran aquellas con unidades más grandes (como
las relacionadas con \verb|area| en nuestro caso) tendrían más peso en el
término y serían las primeras en perder influencia. Si bien el paquete
\{glmnet\} normaliza las variables de forma automática, se optó por dejar por
explícito este paso del preprocesamiento de los datos en el código de R.

Se utilizó la implementación del paquete \verb|{glmnet}| mediante
\verb|{tidymodels}| de R. Este paquete, en realidad, ofrece modelos lineales
generalizados con penalización \textit{elastic net}, que es una combinación
lineal de la penalización LASSO y \textit{Ridge}
\parencite{tay_etal_2023_elastic}. El hiperparámetro que controla esa
combinación es $\alpha$ (\verb|mixture| en la implementación usada), que se
fijó en 1 para utilizar exclusivamente la penalización LASSO. Esto se hizo, por
un lado, para simplificar la búsqueda de hiperparámetros óptimos y, por otro
lado, porque es razonable que en este escenario sea útil seleccionar variables
predictoras para luego priorizar qué aspectos de las células de los tumores
medir. La búsqueda se hizo de forma análoga al modelo de bosques aleatorios; se
creó una grilla de 25 valores equiespaciados en escala logarítmica de base 10
entre $10^{-6}$ y $10^{-1}$ y se seleccionó el que resultase en la mayor área
bajo la curva \textit{precision-recall}. Se utilizaron los mismos
\textit{folds} que en la validación del modelo de bosques aleatorios.

% Métricas de validación de lasso
\begin{figure}[h]
\includegraphics[width=\textwidth]{figuras/lasso1\_metricas.pdf}
\caption{Métricas de desempeño de la regresión logística con penalización LASSO según parámetro de penalización $\lambda$.}
\label{fig:lassometricas}
\end{figure}

La Figura \ref{fig:lassometricas} muestra el valor de las mismas tres métricas
que en el modelo anterior según el valor del hiperparámetro $\lambda$. El mejor
valor encontrado fue $\lambda$ = 0.0056 con un área de 0.991 y error estándar
de 0.0031 aproximadamente. Se ve que cuando crece el coeficiente de
penalización la precisión (con umbral 0.5) sube y el \textit{recall} cae,
seguramente porque el modelo está prediciendo mayormente negativos, lo cual
reduce los falsos positivos y aumenta los falsos negativos. Para que se dé este
resultado, de todas formas, los verdaderos positivos no deberían bajar al mismo
ritmo que los falsos positivos, caso contrario la precisión no debería cambiar
mucho.

% Acumulada de las probabilidades predichas de lasso y umbral vs F2 
\begin{figure}[h]
  \includegraphics[width=\textwidth]{figuras/lasso1\_pred\_umbral.pdf}
  \caption{(A) Distribución acumulada de las probabilidades predichas y (B) puntaje $F_2$ según umbral de decisión correspondientes al modelo de regresión logística con penalización LASSO. La línea punteada vertical indica el umbral de decisión elegido.}
\label{fig:lassopredumb}
\end{figure}

La búsqueda del umbral óptimo se realizó de la misma manera que en el modelo
anterior. La Figura \ref{fig:lassopredumb} es análoga a la figura
\ref{fig:rfpredumb}. Se ve que este modelo también logra una buena separación
de los casos y, en comparación con el anterior, asigna probabilidades más altas
a los tumores malignos, con alrededor del 70\% de las observaciones malignas
con probabilidad cercana a 1. Sin embargo, sigue habiendo casos de este tipo
con predicciones más bajas de lo deseable que se asemejan a las de los casos
benignos con predicciones más altas. El umbral que maximiza el puntaje $F_2$
(calculado sobre el total de predicciones sobre los \textit{folds}) es 0.2 con
un valor $F_2$ de 0.9595.

El modelo final estimó 12 parámetros distintos de cero, incluyendo el
intercepto, que se listan en el Cuadro \ref{fig:lassoespec}, donde se resumen las
especificaciones. Sorprendentemente, la mayor parte de las variables
seleccionadas fueron los peores casos (\verb|_worst|) de ciertas variables,
específicamente, del radio, la textura, el perímetro, la suavidad, la
concavidad, la cantidad de secciones cóncavas y la simetría. Fueron
seleccionadas tres variables de desvío estándar (\verb|_se|): la del radio, la
compacidad y la dimensión fractal. Sólo una corresponde a la media
(\verb|_mean|), la del número de secciones cóncavas. La mayoría de las
variables aumentan las chances de que el tumor sea maligno; sólo el desvío de
la compacidad y de la dimensión fractal reducen esas chances.


\section{Conclusiones}\label{sec:con}

% Sección: Conclusiones

Los resultados del ejercicio fueron positivos; todos los modelos alcanzaron
métricas sumamente altas en el conjunto de testeo con muy pocos casos mal
clasificados y, más importantemente en algunos casos, sin falsos negativos.
Este resultado, sin embargo, probablemente diga más sobre la calidad de los
datos que sobre el ejercicio de modelado. El conjunto de datos no contaba con
valores faltantes, no parecía haber errores en la información, y en general
había fuertes asociaciones entre las variables predictoras y la respuesta.

Uno de los dos mejores modelos es paramétrico y el otro no paramétrico. Bosques
aleatorios tuvo un desempeño marginalmente mejor, pero la regresión logística
con LASSO tiene la ventaja de que, por un lado, permite reducir la cantidad de
variables que es necesario medir para obtener una predicción y, por otro lado,
los coeficientes que estima resultan más interpretables, sobre todo en este
caso que corresponden a las variables originales normalizadas. Es posible
extraer la importancia de cada variable de un modelo de bosques aleatorios,
pero las medidas de importancia no son tan directamente interpretables como los
coeficientes de una regresión, que tienen vínculo directo con el modelo
probabilístico que usamos para conceptualizar el fenómeno. Por estos motivos
considero que es preferible el modelo logístico por sobre el de bosques
aleatorios.

Una de las principales dificultades del trabajo fue la falta de conocimiento
sustantivo sobre el área del conocimiento a la que pertenecen los datos.
Resultó difícil generar ideas a partir de las distribuciones y asociaciones
observadas en la etapa de exploración al no tener noción de la significancia
científica de las variables. Finalmente se utilizaron todas las variables
disponibles en casi todos los modelos ya que no hubo motivos sustantivos para
excluir, combinar o transformarlas más allá de lo que exigían los modelos. Esto
también impacto en la selección de la métrica de evaluación. Se optó por usar
el puntaje $F_\beta$, pero, sin conocimiento del área, resultó prácticamente
imposible decidir una ponderación adecuada. Por otra parte, el error puede
haber estado, justamente, en la elección de esa métrica.


% Bibliografía ---------------------------------------------------------------
\printbibliography
\clearpage

% ----------------------------------------------------------------------------
\appendix

\section{Especificaciones de los modelos seleccionados}\label{sec:anexo1}

% Especificación del modelo lasso
\begin{figure}[h]
  \renewcommand{\figurename}{Cuadro}
\begin{framed}
\begin{itemize}
  \item Implementación: R con el paquete \verb|{glmnet}| mediante \verb|{tidymodels}|
  \item Hiperparámetros
  \begin{itemize}
    \item Coeficiente de penalización $\lambda$: 0.0056
    \item Coeficiente de mixtura $\alpha$: 1 (exclusivamente LASSO)
    \item Umbral de decisión: 0.2
  \end{itemize}
  \item Preprocesamiento: Normalización de todas las variables predictoras
  \item Variables predictoras (entrenamiento): Todas las del conjunto de datos
  \item Variables predictoras (seleccionadas): 
\end{itemize}
  \centering
  \input{cuadros/lasso1\_coefs.tex}
\end{framed}
\caption{Especificaciones del modelo de regresión logística con penalización LASSO}
\label{fig:lassoespec}
\end{figure}

% Especificación del modelo de bosques aleatorios
\begin{figure}[h]
  \renewcommand{\figurename}{Cuadro}
\begin{framed}
\begin{itemize}
  \item Implementación: R con el paquete \verb|{ranger}| mediante \verb|{tidymodels}|
  \item Semilla: 260224
  \item Hiperparámetros
  \begin{itemize}
    \item Árboles por bosque: 100
    \item Variables a probar: 1
    \item Mín. observaciones por hoja: 10
    \item Umbral de decisión: 0.29
  \end{itemize}
  \item Preprocesamiento: Normalización de todas las variables predictoras
  \item Variables predictoras: Todas menos
  \begin{itemize}
    \item \verb|fractal_dimension_mean|
    \item \verb|fractal_dimension_se|
    \item \verb|fractal_dimension_worst|
    \item \verb|texture_se|
    \item \verb|smoothness_se|
    \item \verb|symmetry_se|
  \end{itemize}
\end{itemize}
\end{framed}
\caption{Especificaciones del modelo de bosques aleatorios}
\label{fig:rfespec}
\end{figure}

\clearpage

% Cuadros --------------------------------------------------------------------
\section{Cuadros y gráficos exploratorios}\label{sec:anexo2}

% Medidas resumen de las variables
\begin{table}[h]
\small
\begin{tabular*}{\textwidth}{@{\extracolsep{\fill}}lrrrrr}
\toprule
Característica & Mín. & Máx. & Mediana & Media & s.d. \\
\midrule
  \multicolumn{6}{c}{Variables con sufijo \textit{\_mean}} \\
\midrule
radius & 6.98 & 27.42 & 13.40 & 14.12 & 3.45\\
texture & 9.71 & 39.28 & 18.92 & 19.37 & 4.49\\
perimeter & 43.79 & 186.90 & 86.74 & 91.89 & 23.66\\
area & 143.50 & 2501.00 & 552.95 & 652.74 & 345.16\\
smoothness & 0.05 & 0.16 & 0.10 & 0.10 & 0.01\\
compactness & 0.02 & 0.31 & 0.09 & 0.10 & 0.05\\
concavity & 0.00 & 0.41 & 0.06 & 0.09 & 0.08\\
concave & 0.00 & 0.19 & 0.03 & 0.05 & 0.04\\
symmetry & 0.12 & 0.27 & 0.18 & 0.18 & 0.03\\
fractal\_dimension & 0.05 & 0.10 & 0.06 & 0.06 & 0.01\\
\midrule
  \multicolumn{6}{c}{Variables con sufijo \textit{\_se}} \\
\midrule
radius & 0.11 & 2.55 & 0.31 & 0.39 & 0.26\\
texture & 0.36 & 4.88 & 1.08 & 1.19 & 0.54\\
perimeter & 0.77 & 18.65 & 2.23 & 2.79 & 1.86\\
area & 6.80 & 542.20 & 23.66 & 39.10 & 42.09\\
smoothness & 0.00 & 0.02 & 0.01 & 0.01 & 0.00\\
compactness & 0.00 & 0.14 & 0.02 & 0.03 & 0.02\\
concavity & 0.00 & 0.40 & 0.03 & 0.03 & 0.03\\
concave & 0.00 & 0.05 & 0.01 & 0.01 & 0.01\\
symmetry & 0.01 & 0.06 & 0.02 & 0.02 & 0.01\\
fractal\_dimension & 0.00 & 0.03 & 0.00 & 0.00 & 0.00\\
\midrule
  \multicolumn{6}{c}{Variables con sufijo \textit{\_worst}} \\
\midrule
radius & 7.93 & 36.04 & 14.98 & 16.30 & 4.79\\
texture & 12.02 & 49.54 & 25.56 & 25.91 & 6.38\\
perimeter & 50.41 & 251.20 & 97.98 & 107.45 & 33.12\\
area & 185.20 & 4254.00 & 691.75 & 882.35 & 573.72\\
smoothness & 0.07 & 0.22 & 0.13 & 0.13 & 0.02\\
compactness & 0.03 & 1.06 & 0.22 & 0.26 & 0.15\\
concavity & 0.00 & 1.25 & 0.23 & 0.28 & 0.21\\
concave & 0.00 & 0.29 & 0.10 & 0.11 & 0.06\\
symmetry & 0.16 & 0.66 & 0.28 & 0.29 & 0.06\\
fractal\_dimension & 0.06 & 0.21 & 0.08 & 0.08 & 0.02\\
\bottomrule
\end{tabular*}

  \caption{Estadísticos descriptivos de las variables según sufijo (\textit{\_mean, \_se, \_worst}) y características de las células}
\label{tab:resumen}
\end{table}

% Gráficos -------------------------------------------------------------------
% Gráficos radius
\begin{figure}[h]
\includegraphics[width=\textwidth, center]{figuras/distribuciones\_radius.pdf}
\caption{Histogramas de densidad de las variables asociadas a \textit{radius} según \textit{target} (fila superior) y diagramas de dispersión de las mismas variables de a pares (fila inferior)}
\label{fig:radius}
\end{figure}

% Gráficos texture
\begin{figure}[h]
\includegraphics[width=\textwidth, center]{figuras/distribuciones\_texture.pdf}
\caption{Histogramas de densidad de las variables asociadas a \textit{texture} según \textit{target} (fila superior) y diagramas de dispersión de las mismas variables de a pares (fila inferior)}
\label{fig:texture}
\end{figure}

% Gráficos perimeter
\begin{figure}[h]
\includegraphics[width=\textwidth, center]{figuras/distribuciones\_perimeter.pdf}
\caption{Histogramas de densidad de las variables asociadas a \textit{perimeter} según \textit{target} (fila superior) y diagramas de dispersión de las mismas variables de a pares (fila inferior)}
\label{fig:perimeter}
\end{figure}

%% Gráficos area
%\begin{figure}[h]
%\includegraphics[width=\textwidth, center]{figuras/distribuciones\_area.pdf}
%\caption{Histogramas de densidad de las variables asociadas a \textit{area} según \textit{target} (fila superior) y diagramas de dispersión de las mismas variables de a pares (fila inferior)}
%\label{fig:area}
%\end{figure}

% Gráficos smoothness
%\begin{figure}[h]
%\includegraphics[width=\textwidth, center]{figuras/distribuciones\_smoothness.pdf}
%\caption{Histogramas de densidad de las variables asociadas a \textit{smoothness} según \textit{target} (fila superior) y diagramas de dispersión de las mismas variables de a pares (fila inferior)}
%\label{fig:smooth}
%\end{figure}

% Gráficos compactness
\begin{figure}[h]
\includegraphics[width=\textwidth, center]{figuras/distribuciones\_compactness.pdf}
\caption{Histogramas de densidad de las variables asociadas a \textit{compactness} según \textit{target} (fila superior) y diagramas de dispersión de las mismas variables de a pares (fila inferior)}
\label{fig:compact}
\end{figure}

% Gráficos concavity
\begin{figure}[h]
\includegraphics[width=\textwidth, center]{figuras/distribuciones\_concavity.pdf}
\caption{Histogramas de densidad de las variables asociadas a \textit{concavity} según \textit{target} (fila superior) y diagramas de dispersión de las mismas variables de a pares (fila inferior)}
\label{fig:concavity}
\end{figure}

% Gráficos concave
\begin{figure}[h]
\includegraphics[width=\textwidth, center]{figuras/distribuciones\_concave.pdf}
\caption{Histogramas de densidad de las variables asociadas a \textit{concave} según \textit{target} (fila superior) y diagramas de dispersión de las mismas variables de a pares (fila inferior)}
\label{fig:concave}
\end{figure}

% Gráficos symmetry
\begin{figure}[h]
\includegraphics[width=\textwidth, center]{figuras/distribuciones\_symmetry.pdf}
\caption{Histogramas de densidad de las variables asociadas a \textit{symmetry} según \textit{target} (fila superior) y diagramas de dispersión de las mismas variables de a pares (fila inferior)}
\label{fig:symmetry}
\end{figure}

% Gráficos fractal_dimension
\begin{figure}[h]
\includegraphics[width=\textwidth, center]{figuras/distribuciones\_fractal\_dimension.pdf}
\caption{Histogramas de densidad de las variables asociadas a \textit{fractal\_dimension} según \textit{target} (fila superior) y diagramas de dispersión de las mismas variables de a pares (fila inferior)}
\label{fig:fractal}
\end{figure}

\end{document} % --------------------------------------------------------------

